\section{模型分析检验}

\subsection{主要结论的证据等级}

不同检验支撑的结论强度不同。表\ref{tab:evidence}按第 \ref{sec:evidence-classes} 节的五个等级归类
各问的主要结论，同一问中的直接证据、模型条件推论、半合成校准、外推与敏感性分别表述，不互相替代。

\begin{table}[htbp]
  \centering
  \caption{各问主要结论的证据等级}
  \label{tab:evidence}
  \small
  \setlength{\tabcolsep}{4pt}
  \begin{tabular}{c>{\raggedright\arraybackslash}p{2.72cm}>{\raggedright\arraybackslash}p{2.72cm}%
                  >{\raggedright\arraybackslash}p{2.35cm}>{\raggedright\arraybackslash}p{2.95cm}%
                  >{\raggedright\arraybackslash}p{2.55cm}}
    \toprule
    问题 & A. 直接证据 & B. 模型条件推论 & C. 半合成校准 & D. 外推与敏感性 & E. 诊断与未决 \\
    \midrule
    问题一 & 域级质量评分与冲突率；1M 留出集检验通过 & 1M 响应面上的领域效应 & --- & 11 种方法变体；60M 与 1B 上的迁移检验；A13/A15 压力诊断 & 冲突成因未逐对检验 \\
    \addlinespace
    问题二 & B5、B4 盲验证的中位绝对相对误差 7.52\%、8.07\% & 候选形式下的边际效用、弹性与替代条件 & B7 上的候选形式（未采纳）；B2 对照 & 有效域以外的计算 & B1 的生成公式识别；B8 隔离 \\
    \addlinespace
    问题三 & --- & 各预算、各上下文长度下的最优配置与结构性转移点 & --- & 上下文长度敏感性；删一稳健性；三族原始质量成本比较 & 非基线质量配置待 $Q_0$ 确定 \\
    \addlinespace
    问题四 & 历史前沿与得分分解；算力子集（54 个模型）上的分解；6 个月回测 & 前沿集合的事后分解 & --- & 12/24 个月预测；算力放缓的转移敏感性；时间口径与总体定义 & 损失–得分映射识别不出斜率 \\
    \bottomrule
    \multicolumn{6}{l}{\footnotesize 注：“---”表示该问题在相应证据等级下无对应结论。} \\
  \end{tabular}
\end{table}

\subsection{灵敏度分析}

\textbf{问题一：}11 种预先声明的方法变体下，域级质量评分的最高与最低领域不变，中间领域的
排序随方法变化；缺失记录的三种极端取值使评分变化不超过 0.0002。配比模型的领域效应在 60M 的
配对设计上与 1M 保持 96.2\% 的符号一致，但在规模与设计同时改变的 1B 上只有 68.3\%。

\textbf{问题二：}删去任一条训练轨迹，参数变化不超过估计值的 0.0905\%；剖面目标函数在各参数
偏离 2\% 时都显著上升。二者说明参数由 B1 唯一确定，但由于 B1 与已发表公式逐点一致，这种
稳定性来自数据的生成方式，而不说明真实模型的规律同样确定。

\textbf{问题三：}上下文长度由 2048 增至 131072 时，两个转移点都移向更高的预算，$10^{19}$ FLOPs
下的最优参数量约减半；区间变化在观测取值之间只能被夹住。8 组删一参数下，不受边界约束的最优配置
相对变化不超过 $3.8\times10^{-4}$，每个代表性预算所处的区间不变；预测损失的差异在发布精度下
无法分辨。

\textbf{问题四：}12 个月预测对时间口径最敏感（50.37 至 65.31），其次是总体定义（纳入模型
合并、放宽开源条件、分组）与算力占比的转移区间；时间口径的影响超过模型本身的预测区间宽度。

\subsection{误差分析}

\textbf{问题一：}在同规模留出集上，13 个目标中 Pile-CC 的平均绝对相对误差最小（1.6\%，$R^2=0.865$），
DM Mathematics 最大（34.0\%，$R^2=0.413$），各领域的可预测程度差异明显，宏平均指标会掩盖
这种差异，逐目标结果见图\ref{fig:q1-transfer}。

\textbf{问题二：}按 B1 重新拟合的经典标度律在独立观测数据上的中位绝对相对误差约为 8\%；
它在拟合数据 B1 本身上的样本内中位绝对相对误差为 $3.26\times10^{-5}$，不是验证统计量，也不同于
第 \ref{sec:q2-generator} 节把已发表常数直接代入 B1 得到的 $3.24\times10^{-5}$。前者约 8\% 才是其对真实
模型的预测误差量级。

\textbf{问题三：}最优配置的数值误差来自求解与参数两方面：与不使用解析候选的数值搜索相比，
最优配置的相对差异不超过 $2.65\times10^{-10}$；参数方面的影响见上文的删一稳健性。这两项都不包含
经典标度律本身相对真实模型的误差（独立观测数据上的中位绝对相对误差约 8\%，第 \ref{sec:q2-validity} 节），
因此不能把配置结果的数值精度理解为对真实训练的预测精度。

\textbf{问题四：}6 个月回测中时间趋势模型的平均偏差为 $-5.49$ 分，每次都低估了前沿，说明
线性延续在近期是保守的；正式预测区间因此计入了回测误差。
